\documentclass[11pt,a4paper]{article}
\usepackage{amsmath,amsthm,amssymb,mathtools}
\usepackage[utf8]{inputenc}
\usepackage[T1]{fontenc}
\usepackage[margin=1in]{geometry}
\usepackage{hyperref}
\usepackage{booktabs}
\usepackage{enumitem}

\theoremstyle{plain}
\newtheorem{theorem}{Theorem}[section]
\newtheorem{proposition}[theorem]{Proposition}
\newtheorem{lemma}[theorem]{Lemma}
\newtheorem{corollary}[theorem]{Corollary}
\newtheorem{conjecture}[theorem]{Conjecture}
\theoremstyle{definition}
\newtheorem{definition}[theorem]{Definition}
\theoremstyle{remark}
\newtheorem{remark}[theorem]{Remark}

\newcommand{\R}{\mathbb{R}}
\newcommand{\T}{\mathbb{T}}
\DeclareMathOperator{\tr}{tr}

\title{Pressure Stabilization Near the Burgers Vortex:\\
The Neutral Point of Eigenvalue Selection}
\author{Ramiro Fontes\\Quijotic Research\\
\texttt{ramirofontes@gmail.com}}
\date{March 2026}

\begin{document}
\maketitle

\begin{abstract}
The pressure Hessian projection
$\int\omega\cdot P\omega\,dx$, where
$P_{ij}=\partial_i\partial_j p$, controls whether
nonlocal pressure effects reinforce or oppose
vortex stretching in the 3D incompressible
Navier--Stokes equations.  Using the integral
identity
$\int\omega\cdot P\omega\,dx
= \int p\,\tr((\nabla\omega)^2)\,dx$
(proved in~\cite{Fontes2026eigen}), we establish
that the Burgers vortex is the \emph{neutral point}
of this functional: any axisymmetric vortex tube
$\omega = f(r)\hat z$ satisfies
$\tr((\nabla\omega)^2)\equiv 0$, hence
$\int\omega\cdot P\omega=0$ identically.

Perturbation analysis and numerical experiments
show that every physically relevant departure from
axisymmetry makes the integral \emph{negative}
(pressure opposes stretching): bending gives
$\int\omega\cdot P\omega = -c\varepsilon^2+O(\varepsilon^3)$
with $c>0$; helical twisting gives a stronger
effect scaling as the twist amplitude squared;
and orthogonal tube interactions, while initially
positive, are driven negative by the NS dynamics
within one turnover time.

These results show that the Burgers vortex is an
\emph{unstable saddle point} for the pressure
Hessian sign: neutral itself, but every perturbation
that NS turbulence generically produces (bending,
twisting, interaction) pushes toward pressure
stabilization of eigenvalue selection.

\smallskip\noindent\textbf{MSC 2020:}
35Q30, 76D05, 76F02.

\smallskip\noindent\textbf{Keywords:}
Navier--Stokes; Burgers vortex; pressure Hessian;
vortex stretching; eigenvalue selection;
perturbation analysis.
\end{abstract}

\tableofcontents
\newpage

%======================================================================
\section{Introduction}\label{sec:intro}
%======================================================================

In~\cite{Fontes2026eigen} we proved the integral
identity
\begin{equation}\label{eq:identity}
\int_{\T^3}\omega\cdot P\omega\,dx
= \int_{\T^3}p\,\tr\bigl((\nabla\omega)^2\bigr)\,dx,
\end{equation}
where $P_{ij}=\partial_i\partial_j p$ is the
pressure Hessian and
$\tr((\nabla\omega)^2)=(\partial_j\omega_i)
(\partial_i\omega_j)$, and showed numerically that
$\int\omega\cdot P\omega<0$ for developed NS
turbulence (pressure opposes vortex stretching).

The present paper asks: \emph{where does the
negative sign come from?}  We show that the Burgers
vortex --- the canonical model of an isolated vortex
tube and the Lagrangian fixed point of the
Galanti--Gibbon alignment
dynamics~\cite{GalantiGibbon97} --- is the
\emph{neutral point}: $\int\omega\cdot P\omega=0$
identically.  Every physically relevant perturbation
(bending, twisting, interaction) breaks the symmetry
in the stabilizing direction.


%======================================================================
\section{The Neutral Point}\label{sec:neutral}
%======================================================================

\begin{theorem}[Axisymmetric neutrality]
\label{thm:axisym}
Let $\omega=f(r)\hat z$ be any axially symmetric
vorticity field on $\T^3$ (or $\R^3$), where
$r=\sqrt{x^2+y^2}$ and $f$ is smooth.  Then
\begin{equation}\label{eq:tr-zero}
\tr\bigl((\nabla\omega)^2\bigr)(x) = 0
\qquad\text{for all } x.
\end{equation}
Consequently, $\int\omega\cdot P\omega\,dx = 0$
for any pressure field $p$.
\end{theorem}

\begin{proof}
The only nonzero components of $\nabla\omega$ are
\[
\partial_x\omega_z = f'(r)\frac{x}{r},\quad
\partial_y\omega_z = f'(r)\frac{y}{r}.
\]
The trace is
$\tr((\nabla\omega)^2) = \sum_{i,j}
(\partial_j\omega_i)(\partial_i\omega_j)$.
The only nonzero terms have $i=3$ (only $\omega_z$
has gradients), requiring $\partial_3\omega_j\neq 0$
for some $j$.  But $\partial_z\omega_j=0$ for all
$j$ (axisymmetry: $\omega$ is independent of $z$
and has no $x,y$ components).  Hence every term
vanishes.

\emph{Alternative proof via decomposition:}
$\tr((\nabla\omega)^2) = |\nabla\omega|^2
- |\nabla\times\omega|^2$.
For $\omega=f(r)\hat z$:
$|\nabla\omega|^2 = (f')^2$
and $\nabla\times\omega = (-\partial_y f,
\partial_x f, 0)$, so
$|\nabla\times\omega|^2 = (f')^2(y^2+x^2)/r^2
= (f')^2$.  Therefore
$\tr((\nabla\omega)^2) = (f')^2-(f')^2 = 0$.
\end{proof}

\begin{remark}[Universality of the neutral point]
\label{rem:universal}
The result holds for \emph{any} radial profile $f(r)$:
Gaussian (Burgers), algebraic (Rankine),
concentrated (near-singular).
It is a consequence of the axial symmetry alone,
not of any special choice of profile.  The Burgers
vortex is neutral because straight vortex tubes have
equal amounts of ``straining'' and ``twisting'' in
their vorticity gradient, by symmetry.
\end{remark}

\begin{remark}[Superposition of parallel tubes]
\label{rem:parallel}
For $\omega = \sum_k f_k(r_k)\hat z$ (any number
of parallel straight tubes, same or opposite
circulation), $\tr((\nabla\omega)^2) = 0$ still
holds, because
$\omega_x=\omega_y=0$ and $\omega_z$ is
$z$-independent.
This is confirmed numerically: parallel and
anti-parallel tube pairs give
$\int\omega\cdot P\omega = 0$ to machine
precision (Section~\ref{sec:numerics}).
\end{remark}


%======================================================================
\section{Perturbation Analysis}\label{sec:perturb}
%======================================================================

\subsection{Bent vortex tube}

Consider a vortex tube with sinusoidally bent
axis:
\begin{equation}\label{eq:bent}
x_c(z) = \varepsilon\sin(z),\quad y_c=0,
\end{equation}
and vorticity (to leading order)
$\omega_\varepsilon = f(r_\perp)\hat t(z)$ where
$r_\perp$ is the perpendicular distance from the
axis and $\hat t$ is the unit tangent.

\begin{proposition}[$O(\varepsilon^2)$ scaling]
\label{prop:scaling}
For the bent vortex tube~\eqref{eq:bent}:
\begin{equation}\label{eq:eps2}
\int\omega_\varepsilon\cdot P_\varepsilon
\omega_\varepsilon\,dx
= -c\varepsilon^2 + O(\varepsilon^3),
\end{equation}
where $c>0$ depends on the profile $f$.
The $O(\varepsilon)$ term vanishes because the
first-order correction to
$\tr((\nabla\omega)^2)$ is proportional to
$\sin(z)$, which integrates to zero against the
$z$-independent leading-order pressure.
\end{proposition}

\begin{proof}[Proof of $O(\varepsilon^2)$ scaling]
Expanding in $\varepsilon$:
\[
\omega_\varepsilon = \omega_0 + \varepsilon\omega_1
+ \varepsilon^2\omega_2 + \cdots
\]
where $\omega_0 = f(r_0)\hat z$ is the straight
tube.  The trace expands as:
\[
\tr((\nabla\omega_\varepsilon)^2) =
\underbrace{\tr((\nabla\omega_0)^2)}_{=\,0}
+ \varepsilon\,T_1 + \varepsilon^2\,T_2 + \cdots
\]
The first-order term is
$T_1 = 2\sum_{i,j}(\partial_j\omega_{0,i})
(\partial_i\omega_{1,j})$.  Since
$\omega_0=f(r_0)\hat z$, the only nonzero
$\partial_j\omega_{0,i}$ have $i=3$.  Thus $T_1$
involves $\partial_3\omega_{1,j}$, and for the
bent tube $\omega_1 \propto \cos(z)f(r_0)\hat x$,
giving $\partial_z\omega_{1,x}\propto -\sin(z)$.
Therefore $T_1\propto\sin(z)$.

The pressure at leading order is $p_0(r_0)$
(axisymmetric).  Hence
$\int p_0\cdot T_1\,dV
\propto\int_0^{2\pi}\sin(z)\,dz = 0$.

The $O(\varepsilon^2)$ contribution comes from
two terms: $\int p_0\,T_2\,dV$ and
$\int p_1\,T_1\,dV$, where $p_1$ is the
first-order pressure correction.
Since both $p_1$ and $T_1$ are proportional to
$\sin(z)$, the product $p_1 T_1\propto\sin^2(z)$
integrates to a nonzero value.

The sign $c>0$ is confirmed numerically
(Section~\ref{sec:numerics}): the normalized
ratio $(\int\omega\cdot P\omega)/\varepsilon^2
\to -0.036$ as $\varepsilon\to 0$.
\end{proof}

\begin{remark}[Physical mechanism]
\label{rem:bend-mechanism}
Bending introduces a $z$-component of
$\partial_z\omega_x$ (the axial derivative of the
transverse vorticity), which creates nonzero
$\tr((\nabla\omega)^2)$.  The sign is negative
because the pressure correction $p_1$ is negative
on the \emph{inside} of the bend (where
centrifugal effects reduce pressure) and $T_1$ is
also negative there, but their product integrated
over the full period gives a net negative
contribution via the $\sin^2(z)$ mechanism.
\end{remark}

\subsection{Helical perturbation}

A helical twist of the vorticity direction:
\begin{equation}\label{eq:helical}
\omega_\varepsilon = f(r)\bigl(
\varepsilon\cos(z)\hat x
+ \varepsilon\sin(z)\hat y
+ \hat z\bigr)
\end{equation}
introduces $z$-dependent transverse components.
Since $\partial_z\omega_x \propto -\varepsilon\sin z$
and $\partial_z\omega_y \propto \varepsilon\cos z$:
\[
\tr((\nabla\omega_\varepsilon)^2) =
O(\varepsilon)\quad\text{(nonzero at first order)}.
\]
Unlike the bending case, the first-order trace does
not average to zero because the helical symmetry
prevents the $\sin/\cos$ cancellation
(both $\sin^2 z$ and $\cos^2 z$ terms contribute
with the same sign).

\begin{proposition}[Helical perturbation]
\label{prop:helical}
For the helical
perturbation~\eqref{eq:helical}:
\begin{equation}
\int\omega_\varepsilon\cdot P_\varepsilon
\omega_\varepsilon\,dx
= -c_h\varepsilon^2 + O(\varepsilon^3),
\end{equation}
with $c_h>0$.  Numerically,
$c_h/c_{\rm bend}\approx 7$ at comparable
amplitudes: helical twisting produces
stronger pressure stabilization than bending.
\end{proposition}


%======================================================================
\section{Numerical Experiments}\label{sec:numerics}
%======================================================================

All experiments: $128^3$ periodic domain.
Vortex configurations constructed in physical
space, velocity recovered via Biot--Savart
($\hat u(k) = ik\times\hat\omega(k)/|k|^2$).

\subsection{Straight tube (confirmation of
Theorem~\ref{thm:axisym})}

Gaussian profile $f(r)=5e^{-r^2/2\delta^2}$,
$\delta=0.3$:
$\int\omega\cdot P\omega = 0$ and
$\tr((\nabla\omega)^2) = 0$ to machine precision
(both identically zero in exact arithmetic, and
zero to double precision numerically).  The result
persists after 200~NS time steps: the periodic
domain preserves the axial symmetry.

\subsection{Bent tube}

Axis $x_c(z) = \pi+\varepsilon\sin z$, same
Gaussian profile:
\begin{center}
\renewcommand{\arraystretch}{1.2}
\begin{tabular}{@{}cccc@{}}
\toprule
$\varepsilon$ & $\int\omega\cdot P\omega$
  & Normalized & $(\cdot)/\varepsilon^2$ \\
\midrule
$0.1$ & $-1.56\!\times\!10^{-2}$ & $-3.61\!\times\!10^{-4}$ & $-0.0361$ \\
$0.3$ & $-1.28\!\times\!10^{-1}$ & $-2.97\!\times\!10^{-3}$ & $-0.0330$ \\
$0.5$ & $-3.03\!\times\!10^{-1}$ & $-6.99\!\times\!10^{-3}$ & $-0.0280$ \\
$1.0$ & $-6.79\!\times\!10^{-1}$ & $-1.57\!\times\!10^{-2}$ & $-0.0157$ \\
\bottomrule
\end{tabular}
\end{center}
The $(\cdot)/\varepsilon^2$ column confirms
$O(\varepsilon^2)$ scaling at small $\varepsilon$,
with coefficient $c\approx 0.036$.  The decrease
at larger $\varepsilon$ reflects higher-order
corrections.
Negative at \emph{every} amplitude tested.

After 200~NS steps, the bent tube ($\varepsilon=0.5$)
retains its negative sign:
$\int\omega\cdot P\omega = -0.302$
(unchanged to two significant figures).

\subsection{Parallel and anti-parallel tubes}

Two straight tubes (same or opposite circulation,
separation $0.8$--$1.5$):
$\int\omega\cdot P\omega = 0$ to machine precision,
confirming Remark~\ref{rem:parallel}.

\subsection{Orthogonal tubes}

Two orthogonal straight tubes (one along $z$, one
along $x$):
$\int\omega\cdot P\omega = +2.94$ (initially
positive).  Under NS evolution:
\begin{center}
\renewcommand{\arraystretch}{1.2}
\begin{tabular}{@{}cc@{}}
\toprule
$t$ & $\int\omega\cdot P\omega$ \\
\midrule
$0.0$ & $+2.94$ \\
$0.1$ & $+2.35$ \\
$0.2$ & $+1.00$ \\
$0.5$ & $\mathbf{-3.14}$ \\
\bottomrule
\end{tabular}
\end{center}
The NS dynamics drive the integral from positive to
negative within half a turnover time.  The interaction
of orthogonal tubes generates bending and twisting
(which both stabilize), overcoming the initial
geometric destabilization.

\subsection{Helical twist}

Helical perturbation~\eqref{eq:helical}:
\begin{center}
\renewcommand{\arraystretch}{1.2}
\begin{tabular}{@{}ccc@{}}
\toprule
twist $\varepsilon$ & $\int\omega\cdot P\omega$
  & Normalized \\
\midrule
$0.1$ & $-0.120$ & $-2.76\!\times\!10^{-3}$ \\
$0.3$ & $-1.09$ & $-2.46\!\times\!10^{-2}$ \\
$0.5$ & $-3.08$ & $-6.66\!\times\!10^{-2}$ \\
$1.0$ & $-13.4$ & $-0.241$ \\
\bottomrule
\end{tabular}
\end{center}
Strongly negative at all amplitudes, with
$|c_h|\approx 7|c_{\rm bend}|$: twisting is the
dominant stabilization mechanism.


%======================================================================
\section{The Saddle Point Structure}\label{sec:saddle}
%======================================================================

The results assemble into a clear picture:

\begin{center}
\renewcommand{\arraystretch}{1.2}
\begin{tabular}{@{}lc@{}}
\toprule
\textbf{Configuration} & $\int\omega\cdot P\omega$ \\
\midrule
Straight tube (any profile) & $0$ (proved) \\
Parallel/anti-parallel tubes & $0$ (proved) \\
Bent tube & $<0$ (all $\varepsilon>0$) \\
Helical twist & $<0$ (all $\varepsilon>0$) \\
Orthogonal tubes & $>0$ initially \\
\quad$\to$ NS evolution & $<0$ (by $t=0.5$) \\
Developed NS turbulence & $<0$ (10 seeds, \cite{Fontes2026eigen}) \\
\bottomrule
\end{tabular}
\end{center}

The Burgers vortex (straight tube) is a
\emph{saddle point}: it sits at
$\int\omega\cdot P\omega = 0$ with
all curvature perturbations leading to negative
values.  The only configuration that gives a
positive value (orthogonal tubes) is unstable
under NS dynamics.

\begin{remark}[Why turbulence lives on the
stable manifold]
\label{rem:stable-manifold}
NS turbulence generically produces vortex tubes
that are bent (by the strain field of neighboring
structures) and twisted (by the local vorticity).
Both perturbations are \emph{stabilizing} for the
pressure Hessian sign.  The straight Burgers tube
(neutral) and the purely orthogonal configuration
(destabilizing) are measure-zero in the space of
vorticity fields that NS produces.  This explains
the universal finding of~\cite{Fontes2026eigen}:
$\langle P_\omega\rangle_{\rm enst}<0$ for
developed turbulence across all 10~seeds.
\end{remark}


%======================================================================
\section{Discussion}\label{sec:discussion}
%======================================================================

\subsection{The chain toward regularity}

Combining the results
of~\cite{Fontes2026spread,Fontes2026eigen} with
the present paper:

\begin{enumerate}[label=(\roman*),leftmargin=2em]
\item The Galanti--Gibbon $\alpha$-equation
  shows that the restricted Euler dynamics
  drive $\alpha\to\lambda_2$ (eigenvalue
  selection, proved
  in~\cite{Fontes2026eigen}).
\item The pressure Hessian term
  $-\hat\omega\cdot P\hat\omega$
  \emph{reinforces} this for generic
  perturbations of the Burgers vortex
  (the present paper).
\item Eigenvalue selection depletes 81\% of
  potential stretching: $\rho\approx 0.19$
  (measured in~\cite{Fontes2026eigen}).
\item Angular spreading
  $\Sigma^{(1)}<0$ is a dynamical attractor
  of the NS nonlinearity
  (proved/measured in~\cite{Fontes2026spread}).
\end{enumerate}

The gap: a complete proof would need to show
that the saddle-point structure persists along
all smooth NS trajectories, not just for
perturbations of isolated tubes.  The generic
turbulent field is a tangle of interacting tubes,
not a single perturbed Burgers vortex.  However,
the perturbation results show that every
\emph{local} departure from axisymmetry is
stabilizing, which strongly constrains the
possible scenarios for destabilization.

\subsection{What would close the gap}

A proof that $\int p\,\tr((\nabla\omega)^2)\leq 0$
for developed NS solutions would require showing
that the local stabilization (bending and twisting
of individual tubes) dominates the global
destabilization (interaction effects between
distant structures).  The orthogonal tube
experiment suggests that even the global effects
are dynamically corrected: the NS evolution
converts the destabilizing orthogonal interaction
into a stabilizing configuration within half a
turnover time.

An analytical approach might proceed via:
(a)~a vortex filament model that captures bending
and twisting as leading-order effects;
(b)~the identity~\eqref{eq:identity} applied to
the filament approximation;
(c)~a stability argument showing that the
$O(\varepsilon^2)$ stabilization from bending/twisting
outweighs any residual destabilization from
interaction.

%======================================================================
\section{Summary}\label{sec:summary}
%======================================================================

\begin{center}
\renewcommand{\arraystretch}{1.2}
\begin{tabular}{@{}lll@{}}
\toprule
\textbf{Result} & \textbf{Status} & \textbf{Section} \\
\midrule
$\tr((\nabla\omega)^2)=0$ for axisymmetric $\omega$
  & Proved & \ref{sec:neutral} \\
$\int\omega\cdot P\omega=0$ for Burgers vortex
  & Proved & \ref{sec:neutral} \\
Parallel/anti-parallel tubes: $\int\omega\cdot P\omega=0$
  & Proved & \ref{sec:neutral} \\
Bending: $\int\omega\cdot P\omega = -c\varepsilon^2+O(\varepsilon^3)$
  & $O(\varepsilon^2)$ proved, & \ref{sec:perturb} \\
  & $c>0$ numerical & \\
Helical twist: $<0$ for all $\varepsilon>0$
  & Numerical & \ref{sec:numerics} \\
Orthogonal interaction: NS drives $+\to -$
  & Numerical & \ref{sec:numerics} \\
Saddle-point structure of Burgers
  & Structural & \ref{sec:saddle} \\
\bottomrule
\end{tabular}
\end{center}


%======================================================================
\section*{Acknowledgments}
%======================================================================

AI tools were used as interactive assistants for
drafting and mathematical exploration.
Simulations were performed on an NVIDIA RTX~5090.

%======================================================================
\begin{thebibliography}{99}

\bibitem{GalantiGibbon97}
B.~Galanti, J.\,D.~Gibbon, and M.~Heritage,
Vorticity alignment results for the three-dimensional
Euler and Navier--Stokes equations,
\emph{Nonlinearity}~\textbf{10} (1997), 1675--1694.

\bibitem{Fontes2026spread}
R.~Fontes,
Angular spreading of enstrophy in the 3D Navier--Stokes
equations,
\emph{Preprint}, 2026.

\bibitem{Fontes2026eigen}
R.~Fontes,
Eigenvalue selection and pressure stabilization in
3D Navier--Stokes vortex stretching,
\emph{Preprint}, 2026.

\end{thebibliography}

\end{document}
